\section{Simulation Model}
\label{sec:simulation_model}

The  purpuse of the simulation model is to describe how a stored spin wave is converted into an optical field during retrieval and how this conversion is degraded when atoms move, acquire different phases, or leave the interaction region. The full Maxwell--Bloch writing dynamics are not solved. The stored ground-state coherence is used as the initial condition, and the retrieved field is calculated from the coherent radiation of the atoms that remain coupled to the retrievable collective mode.

Each atom contributing to readout is approximated as an ideal radiating dipole with a common internal orientation. Under this approximation, differences between atoms enter through their positions, phases, and weights, while the single-atom radiation pattern is treated as a common factor. Following the standard element-factor and array-factor decomposition used in antenna theory~\cite{AntenaCourseNotes}, the far-field amplitude in the observation direction $\hat{\mathbf n}$ may be written as
\begin{equation}
    \mathcal{E}_{\mathrm{out}}(\hat{\mathbf n},t)
    \propto
    \mathcal{E}_{\mathrm{dip}}(\hat{\mathbf n})
    \,
    \mathcal{A}(\hat{\mathbf n},t),
    \label{eq:field_element_array_factor}
\end{equation}
where $\mathcal{E}_{\mathrm{out}}(\hat{\mathbf n},t)$ is the emitted field amplitude, $\mathcal{E}_{\mathrm{dip}}(\hat{\mathbf n})$ is the single-dipole radiation factor, and $\mathcal{A}(\hat{\mathbf n},t)$ is the collective array factor. Since the relevant quantity is the relative retrieval into a selected optical mode, $\mathcal{E}_{\mathrm{dip}}(\hat{\mathbf n})$ is taken to be identical for all atoms and is not used to distinguish different trajectories.

The collective part of the field is represented by the weighted coherent sum
\begin{equation}
    \mathcal{A}(\hat{\mathbf n},t)
    =
    \sum_{j=1}^{N}
    W_j(t)
    \exp\left[
        i\Phi_j(\hat{\mathbf n},t)
    \right],
    \label{eq:array_factor_general}
\end{equation}
where $N$ is the number of simulated atoms, $W_j(t)$ is the complex amplitude weight assigned to atom $j$, and $\Phi_j(\hat{\mathbf n},t)$ is the phase with which that atom contributes to the field observed along $\hat{\mathbf n}$. The optical power associated with the selected direction is proportional to the squared modulus of the coherent amplitude,
\begin{equation}
    P(\hat{\mathbf n},t)
    \propto
    \left|
    \mathcal{A}(\hat{\mathbf n},t)
    \right|^2 .
    \label{eq:array_factor_power}
\end{equation}
Constructive collective emission is obtained when the phases in Eq.~\eqref{eq:array_factor_general} remain aligned. Random phases reduce the coherent sum even if individual atoms still radiate, because the field amplitudes no longer add efficiently before the modulus square is taken~\cite{AntenaCourseNotes,PsiL-odamIFm_GorshkovEtAl_2007}.

For spin-wave retrieval, the phase of each atomic contribution is set by the stored spin-wave grating, the read control field, the emitted optical mode, and any internal phase accumulated during storage. Using $\mathbf{k}_{\mathrm{out}}(\hat{\mathbf n})=k_{\mathrm{out}}\hat{\mathbf n}$, the phase is written as
\begin{equation}
    \Phi_j(\hat{\mathbf n},t)
    =
    \left(
        \Delta\mathbf{k}
        +
        \mathbf{k}_{\mathrm{c}}^{\mathrm{read}}
        -
        \mathbf{k}_{\mathrm{out}}(\hat{\mathbf n})
    \right)
    \cdot
    \mathbf{r}_j(t)
    -
    \phi_j(t),
    \label{eq:atom_retrieval_phase}
\end{equation}
where $\Delta\mathbf{k}$ is the spin-wave wave vector written during storage, $\mathbf{k}_{\mathrm{c}}^{\mathrm{read}}$ is the read-control wave vector, $\mathbf{k}_{\mathrm{out}}(\hat{\mathbf n})$ is the retrieved optical wave vector in the observation direction, $\mathbf{r}_j(t)$ is the position of atom $j$ after storage time $t$, and $\phi_j(t)$ is the additional phase accumulated by the ground-state coherence. Constructive retrieval is obtained when the position-dependent phase in Eq.~\eqref{eq:atom_retrieval_phase} is stationary across the ensemble. Following the usual spin-wave phase-matching condition~\cite{PsiL-odamIFm_GorshkovEtAl_2007}, this gives
\begin{equation}
    \mathbf{k}_{\mathrm{out}}
    =
    \Delta\mathbf{k}
    +
    \mathbf{k}_{\mathrm{c}}^{\mathrm{read}}.
    \label{eq:simulation_phase_matching}
\end{equation}
The preferred retrieval direction is therefore fixed by the spatial phase stored in the atomic coherence and by the geometry of the read control field.

The weights $W_j(t)$ describe the amplitude with which each atom contributes to the selected mode. Uniform weights would be obtained for an ideal ensemble illuminated homogeneously by the relevant optical fields. In a finite memory system with nonuniform optical mode profiles, atoms near the centre of the signal and control beams couple more strongly than atoms in the wings. The stored spin wave therefore carries an amplitude envelope, and the coherent sum in Eq.~\eqref{eq:array_factor_general} must be evaluated as a weighted sum.

For a Gaussian optical mode, the transverse field envelope is taken as
\begin{equation}
    u_{\ell}(\mathbf{r})
    =
    \exp\left[
        -\frac{\rho_{\ell}^2(\mathbf{r})}{w_{\ell}^2}
    \right],
    \label{eq:gaussian_field_envelope}
\end{equation}
where the subscript $\ell$ labels the optical field, $\rho_{\ell}(\mathbf{r})$ is the transverse distance from the propagation axis of that field, and $w_{\ell}$ is the Gaussian field radius. The corresponding intensity profile is
\begin{equation}
    I_{\ell}(\mathbf{r})
    \propto
    \left|
    u_{\ell}(\mathbf{r})
    \right|^2
    =
    \exp\left[
        -\frac{2\rho_{\ell}^2(\mathbf{r})}{w_{\ell}^2}
    \right].
    \label{eq:gaussian_intensity_profile}
\end{equation}
If the beam diameter is specified by the full width at half maximum of the intensity, the field radius in Eq.~\eqref{eq:gaussian_field_envelope} is related to $\mathrm{FWHM}_{\ell}$ by
\begin{equation}
    \mathrm{FWHM}_{\ell}
    =
    w_{\ell}\sqrt{2\ln 2}.
    \label{eq:fwhm_to_waist}
\end{equation}

The initial spin-wave profile is determined by the overlap of the signal and write-control fields. Following the spin-wave description of ensemble memories~\cite{PsiL-odamIFm_GorshkovEtAl_2007}, the stored coherence may be written as
\begin{equation}
    S(\mathbf{r},0)
    =
    \mathcal{N}
    \,
    u_{\mathrm{sig}}(\mathbf{r})
    u_{\mathrm{c}}^{*}(\mathbf{r})
    \exp\left[
        i
        \left(
            \mathbf{k}_{\mathrm{sig}}^{\mathrm{write}}
            -
            \mathbf{k}_{\mathrm{c}}^{\mathrm{write}}
        \right)
        \cdot
        \mathbf{r}
    \right],
    \label{eq:initial_spin_wave_profile}
\end{equation}
where $S(\mathbf{r},0)$ is the initial ground-state coherence, $\mathcal{N}$ is a normalization constant, $u_{\mathrm{sig}}(\mathbf{r})$ is the signal-mode envelope, $u_{\mathrm{c}}(\mathbf{r})$ is the write-control envelope, $\mathbf{k}_{\mathrm{sig}}^{\mathrm{write}}$ is the signal wave vector, and $\mathbf{k}_{\mathrm{c}}^{\mathrm{write}}$ is the write-control wave vector. The conjugation of $u_{\mathrm{c}}(\mathbf{r})$ reflects that the stored coherence is written by the signal field relative to the control field.

In the particle representation, the continuous spin wave is sampled by atoms at positions $\mathbf{r}_j(0)$. The initial amplitude weight is
\begin{equation}
    W_j(0)
    =
    \mathcal{N}
    \,
    u_{\mathrm{sig}}\!\left(\mathbf{r}_j(0)\right)
    u_{\mathrm{c}}^{*}\!\left(\mathbf{r}_j(0)\right),
    \label{eq:initial_atomic_weight}
\end{equation}
where $W_j(0)$ is the initial contribution of atom $j$ to the stored mode. The spin-wave phase is kept outside the weight and is included explicitly in Eq.~\eqref{eq:atom_retrieval_phase}. This separation assigns the mode amplitude to $W_j(t)$ and the retrieval interference to $\Phi_j(\hat{\mathbf n},t)$.

Atomic motion is included by evolving the positions $\mathbf{r}_j(t)$ during storage. In a warm vapour, atoms move through the optical mode, collide with buffer-gas atoms, and may encounter the cell walls. Even when the internal coherence survives, the atom contributes from a different position at readout. Since the phase in Eq.~\eqref{eq:atom_retrieval_phase} is evaluated at $\mathbf{r}_j(t)$, motion changes the collective interference pattern and can reduce the phase-matched retrieved field.

For diffusive motion, the stochastic displacement during a time step $\Delta t$ is sampled as~\cite{MonteCarloNotes}
\begin{equation}
    \mathbf{r}_j(t+\Delta t)
    =
    \mathbf{r}_j(t)
    +
    \sqrt{2D\Delta t}\,
    \boldsymbol{\xi}_j,
    \label{eq:diffusive_position_update}
\end{equation}
where $D$ is the diffusion coefficient and $\boldsymbol{\xi}_j$ is a three-dimensional random vector with independent normally distributed components of zero mean and unit variance. This update gives the three-dimensional diffusion result~\cite{MonteCarloNotes}
\begin{equation}
    \left\langle
    \left|
    \mathbf{r}_j(t+\Delta t)-\mathbf{r}_j(t)
    \right|^2
    \right\rangle
    =
    6D\Delta t,
    \label{eq:diffusion_msd}
\end{equation}
where the angle brackets denote an average over trajectories. Other transport models can be incorporated by replacing the update rule for $\mathbf{r}_j(t)$.

A position-dependent ground-state splitting produces an additional phase. The spatially uniform part of the splitting produces only a common phase and is removed from the retrieval amplitude. The relevant contribution is the differential shift $\Delta\omega_{sg}(\mathbf{r})$, which gives
\begin{equation}
    \phi_j(t)
    =
    \int_0^t
    \Delta\omega_{sg}\!\left(\mathbf{r}_j(t')\right)
    \,dt',
    \label{eq:internal_phase_integral}
\end{equation}
where $t'$ is the integration variable. A homogeneous shift gives the same $\phi_j(t)$ for every atom and leaves the collective interference unchanged. Spatially varying shifts give trajectory-dependent phases and reduce the phase-matched spin-wave component.

Wall relaxation and other loss mechanisms are included through a survival factor. The time-dependent atomic weight is written as
\begin{equation}
    W_j(t)
    =
    W_j(0)
    \,
    \eta_j(t)
    \,
    u_{\mathrm{read}}\!\left(\mathbf{r}_j(t)\right),
    \label{eq:time_dependent_weight}
\end{equation}

where $\eta_j(t)$ is the survival factor of the ground-state coherence of atom $j$, and $u_{\mathrm{read}}(\mathbf{r}_j(t))$ is the spatial envelope associated with the readout mode at the atom position. Ideal coherent motion is recovered by setting $\eta_j(t)=1$. Depolarizing events reduce $\eta_j(t)$ and therefore remove part of the atomic contribution from the coherent sum.

Wall-induced relaxation is represented phenomenologically through the mean number of wall collisions that an atom can undergo before losing its ground-state coherence. This parameter, denoted by $N_{\mathrm{b}}$, does not describe the microscopic atom--surface interaction of a specific coating. It provides an effective way to compare uncoated cells with anti-relaxation coatings by changing the probability that coherence survives repeated wall encounters. If atom $j$ has experienced $n_j(t)$ wall collisions by time $t$, the survival factor is taken to be
\begin{equation}
    \eta_j(t)
    =
    \left(
        1-\frac{1}{N_{\mathrm{b}}}
    \right)^{n_j(t)}.
    \label{eq:wall_survival_factor}
\end{equation}
This expression assigns the same survival probability to each wall collision and assumes independent relaxation events. The physical interpretation of $N_{\mathrm{b}}$ and its relation to coated warm-vapour cells will be discuss in section~\ref{sec:coatings}

The simulated retrieval amplitude in a target direction $\hat{\mathbf n}_{\mathrm{tar}}$ is then
\begin{equation}
    \mathcal{A}_{\mathrm{tar}}(t)
    =
    \sum_{j=1}^{N}
    W_j(t)
    \exp\left[
        i\Phi_j(\hat{\mathbf n}_{\mathrm{tar}},t)
    \right],
    \label{eq:target_mode_amplitude}
\end{equation}
where $\mathcal{A}_{\mathrm{tar}}(t)$ is the complex field amplitude emitted into the target optical direction. The corresponding target-mode power is
\begin{equation}
    P_{\mathrm{tar}}(t)
    \propto
    \left|
    \mathcal{A}_{\mathrm{tar}}(t)
    \right|^2 .
    \label{eq:target_mode_power}
\end{equation}

As the detected mode is a fibre mode rather than a single far-field direction, the retrieved amplitude is projected onto the angular fibre mode. The fibre-coupled amplitude is
\begin{equation}
    \mathcal{A}_{\mathrm{fib}}(t)
    =
    \int d\Omega\,
    u_{\mathrm{fib}}^{*}(\hat{\mathbf n})
    \mathcal{E}_{\mathrm{out}}(\hat{\mathbf n},t),
    \label{eq:fibre_overlap_integral}
\end{equation}
where $u_{\mathrm{fib}}(\hat{\mathbf n})$ is the normalized angular mode of the fibre and $d\Omega$ is the differential solid angle. The fibre-coupled power is
\begin{equation}
    P_{\mathrm{fib}}(t)
    \propto
    \left|
    \mathcal{A}_{\mathrm{fib}}(t)
    \right|^2 .
    \label{eq:fibre_power}
\end{equation}
The useful retrieved signal is therefore determined by a coherent mode overlap, rather than by the total optical power emitted into all directions.

The Monte Carlo calculation evaluates the coherent sums by sampling atoms and trajectories. Each sampled atom is assigned an initial position, an optical weight, a stochastic trajectory, a survival factor, and an accumulated internal phase. At each timestep, the retrieval amplitude is evaluated from the current positions, phases, and weights. The numerical result therefore estimates an ensemble average over possible atomic configurations and histories~\cite{MonteCarloNotes}.

For the target direction, the Monte Carlo estimator of the amplitude is
\begin{equation}
    \widehat{\mathcal{A}}_{\mathrm{tar}}(t)
    =
    \frac{1}{N_{\mathrm{MC}}}
    \sum_{j=1}^{N_{\mathrm{MC}}}
    W_j(t)
    \exp\left[
        i\Phi_j(\hat{\mathbf n}_{\mathrm{tar}},t)
    \right],
    \label{eq:monte_carlo_amplitude_estimator}
\end{equation}
where $N_{\mathrm{MC}}$ is the number of Monte Carlo samples and the hat denotes a numerical estimator. The associated power estimator is
\begin{equation}
    \widehat{P}_{\mathrm{tar}}(t)
    \propto
    \left|
    \widehat{\mathcal{A}}_{\mathrm{tar}}(t)
    \right|^2 .
    \label{eq:monte_carlo_power_estimator}
\end{equation}
Under independent sampling, the statistical error of the amplitude estimator scales as $N_{\mathrm{MC}}^{-1/2}$~\cite{MonteCarloNotes}.

The normalized retrieval observable is defined as
\begin{equation}
    \mathcal{R}_{\mathrm{fib}}(t)
    =
    \frac{P_{\mathrm{fib}}(t)}
    {P_{\mathrm{tot}}(0)},
    \label{eq:normalized_fibre_retrieval}
\end{equation}
where $P_{\mathrm{fib}}(t)$ is the fibre-coupled retrieved power after storage time $t$, and $P_{\mathrm{tot}}(0)$ is the total emitted power at zero storage time. This normalization removes the arbitrary absolute scale of the emitted field while retaining the decay of the retrievable mode. The observable is therefore sensitive to phase mismatch, loss of beam overlap, and survival of the ground-state coherence.
